\chapter{Stand des Wissens} \label{cha:standDesWissens}

Die vorliegende Masterarbeit untersucht den Einfluss von Large Language Models (LLMs) auf das Lernen und Prüfen von Programmieranfänger*innen. Der Forschungsstand umfasst mehrere miteinander verknüpfte Bereiche: die besonderen Herausforderungen des Programmierlernens, den Einsatz von LLMs in der Programmierausbildung, deren Rolle in Prüfungs- und Bewertungsszenarien sowie didaktische Faktoren. In diesem Kapitel soll die einschlägige Literatur zusammengeführt werden.

\section{Herausforderungen beim Programmierenlernen im tertiären Bildungsbereich}
Das Erlernen von Programmierkompetenzen gilt als besonders herausfordernd. Mehrere Studien, darunter die Untersuchungen von \textcite{WatsonLi2014FailureRates} sowie \textcite{Simon2019PassRates}, zeigen, dass ein Teil der Studierenden Einführungskurse in die Programmierung nicht erfolgreich abschließt. Die durchschnittlichen Bestehensquoten liegen in den von \textcite{WatsonLi2014FailureRates} untersuchten Kursen bei 67,7~\% und bei \textcite{Simon2019PassRates} bei rund 75~\%.

\textcite[S.~6]{SystematicLiteratureReview2019} identifizieren in ihrer systematischen Literaturübersicht zentrale Einflussfaktoren, die den Erfolg in Einsteigerkursen bestimmen, darunter mathematische Kompetenzen, Problemlösefähigkeiten, Abstraktionsvermögen sowie die Fähigkeit, Fehlermeldungen zu interpretieren und systematisch zu debuggen. Weitere Arbeiten, etwa von \textcite[S.~2]{TeachingIntroductoryQuantitive} oder \textcite[S.~1]{studentPerformance}, zeigen, dass insbesondere die ersten Lernschritte also das Strukturieren eines Problems und der Übergang von natürlicher Sprache zu formalen Anweisungen sowie der Umgang mit Syntax als besonders frustrierend erlebt werden. \textcite[S.~5-7]{Sorva2013NotionalMachines} beschreibt Fehlvorstellungen Programmieranfänger*innen zu Variablen, Wertzuweisungen, Kontrollfluss und Parameterübergabe und wie einzelne Anweisungen den Programmzustand verändern.

Des weiteren werden Herausforderungen bei der Entwicklung eigener Lösungen sowie auch beim Verständnis des vorhandenen Programmcodes verortet. In einer im Jahr 2005 von \textcite[S.~15-16]{Lahtinen2005Difficulties} durchgeführten Befragung von Studierenden und Lehrenden im Bereich der Programmiersprachen C++ und Java wurde foglende Tätigkeiten als besonders schwierig wahrgenommen: Der Entwurf eines Programms für eine vorgegebene Aufgabe, die Aufteilung der Funktionalität in Teilfunktionen sowie das Auffinden von Fehlern (debugging) im eigenen Code. Zudem wurde auch nach Programmier-Konzepten gefragt und wie schwer diese zu erlernen war. Hierbei werden vonseiten der Studierenden einerseits Variablen, Verzweigungen (if und else) und Schleifen als einfach klassifiziert. Andererseits werden Rekursion, Pointer und Referenzen sowie Fehlerbehandlung (error handling) von den Befragten Studierenden als schwer eingeordnet. Ein interessanter Aspekt aus der Studie von \textcite[S.~16]{Lahtinen2005Difficulties} stellt die Frage nach dem Wissenszuwachs dar. Vorlesungen bilden hierbei bei Studierenden und Lehrenden das Schlusslicht, wobei praktische Übungen im Unterricht einen ungemein größeren gefühlten Lernerfolg aufweisen. Am höchsten wird von Studierenden das alleinige Arbeiten and Programmieraufgabaen angegeben, bei den Lehrenden sind es die genannten praktischen Übungen im Unterricht.

Zu den Schwierigkeiten auf Studierenden-Seite zählt auch die Fähigkeit die Programmausführung schrittweise nachvollziehen zu können. \textcite[S.~119-121]{lister_multi-national_2004} zeigen in einer Untersuchung an Hochschulen, dass viele Studierende Probleme haben, die Ergebnisse kurzer Codeabschnitte vorherzusagen und unvollständigen Code passend zu ergänzen. Diese beiden Sachverhalte beziehen sich direkt auf das Konzept des Computational Thinkings (CT) und zeigen einen Mangel der Fähigkeiten auf, die CT unter anderem umschließt. Für CT gibt es unterschiedliche Ansätze, ursprünglich umfasst es nach \textcite[S.~33-34]{wing_computational_2006} die Problemlösung, das Entwerfen von Systemen und auch das Verstehen menschlichen Verhaltens aus dem Blickwinkel der Informatik. In der Literaturübersicht von \textcite{shute_demystifying_2017} werden aus verschiedenen CT-Ansätzen die Kern-Dimensionen Problemzerlegung, Abstraktion, Algorithmenentwurf, Debugging, Iteration und Generalisierung abgeleitet. Das schrittweise Nachvollziehen von Programmcode ist in diesem Zusammenhang eine konkrete Tätigkeit beim Analysieren und Überprüfen eines Algorithmus.

Auf curricularer Ebene bestehen zusätzliche Herausforderungen: Lehrende müssen mit unterschiedlichen Vorkenntnisse umgehen, sinnvolle Aufgabenformate entwickeln und auf eine angemessene Balance zwischen Theorie, Codepraxis und Feedback achten. \textcite[S.~7]{SystematicReviewFinnland}, \textcite[S.~7-8]{DesigningAndDeploying} und \textcite[S.~15-16]{Lahtinen2005Difficulties} zeigen, dass die Wahl der Programmiersprache und der didaktischen Methode (z. B. pair programming, projektorientiertes Lernen, tutorielle Unterstützung) entscheidend dazu beiträgt, wie gut Studierende grundlegende Konzepte wie Kontrollstrukturen, Datenstrukturen oder Objektorientierung verstehen.


\section{Die Verwendung von LLMs in der Programmierausbildung}
Mit dem Aufkommen von generativen Modellen von OpenAI, Google und später auch Anthropic hat sich das Programmierlernen grundlegend verändert. LLMs können Code generieren, erklären, korrigieren und optimieren. Sie können Fehlermeldungen interpretieren, Beispiele erzeugen und Lernende durch ganze Problemlösungsprozesse begleiten. \textcite[S.~1]{DennyEra} beschreiben diesen Wandel als einen Paradigmenwechseln, der sowohl Chancen für individualisiertes Lernen als auch Herausforderungen für Didaktik und Prüfungsdesign mit sich bringt.

Eine Reihe von Studien befasst sich mit den unmittelbaren Lerneffekten. \textcite[S.~1]{RobotsAreComming} sowie \textcite[S.~1]{KazemitabaarCodeAid} zeigen, dass Studierende mit LLM-Unterstützung typischerweise schneller zu funktionierenden Lösungen gelangen und weniger Frustrationspunkte erleben -- insbesondere in frühen Lernphasen. \textcite[S.~565-56]{LeinonenErrorMessages} spricht vom Unterstützungspotential der LLMs und über ihre Fähigkeit, Fehlermeldungen zu erklären und Korrekturvorschläge aufzuzeigen und somit zur Reduktion der Debugging-Zeiten beitragen. \textcite{wang_large_2024} berichten, dass Studierende mit ki-angereicherten Fehlermeldungen weniger zusätzliche Ausführungsversuche zur Fehlerbehebung benötigten als mit Standardmeldungen. Die Studierenden mussten mit Python verschiedene Probleme lösen und wenn sie beim Ausführen des Programms z.B. einen Syntax-Error erzeugt haben, bekammen sie in der Programmierumgebung die für Python tpyische Fehlerbeschreibung und zusätlich auch eine mit GPT erzeugte Erklärung des vorliegenden Fehlers \parencite[S.~1395-1397]{wang_large_2024}.

Auch \textcite[S.~S.1]{santos_its_2025} stellen auch kürzere Fehlerbehebungszeiten im Debuggingprozess mit einer KI-Integration fest. In ihrer Untersuchung stellten sie den Lernenden verschiedene Arten von Fehlerbeschreibungen zur Verfügung: Python nativ, GPT-4o Baseline und GPT-4o fine-tuned. Mit den Fehler-Erklärungen von letzterer konnten die Studierenden Probleme schneller lösen. Jedoch bevorzugen sie die ausführlicheren Beschreibungen vom GPT-4o Basline Modell als die fine-tuned Variante oder die Python nativen Fehlermeldungen.

Gleichzeitig wird in der Forschung intensiv diskutiert, ob die Nutzung solcher Systeme zu oberflächlichen Lernprozessen führt. \textcite[S.~1]{BeckerProgrammingHard} argumentieren, dass Programmieren möglicherweise nicht mehr so schwer ist wie früher, weil LLMs viele kognitive Hürden übernehmen. Damit stellt sich die Frage, welche Kompetenzen Studierende tatsächlich erwerben, wenn wesentliche Teile der Problemlösung ausgelagert werden. Hinweise auf damit verbundene Schwierigkeiten liefert \textcite[S.~484]{prather_widening_2024} in einer Untersuchung mit 21 Programmieranfänger*innen über 21 Einheiten hinweg. Hierbei wurde mittels Interviews, Eye-Tracking und Einzelbeobachtung untersucht, wie die Anfänger*innen mit KI-Tools programmieren. Bei einem Teil der Lernenden, die Probleme bei der Bearbeitung der Aufgaben hatten, beobachteten die Autor*innen, dass die Unterstützung durch KI Fortschritt suggerierte, obwohl weiterhin Verständnisprobleme bestanden. Die Studierenden unterschätzten dabei durch die erfolgreiche Bearbeitung der Aufgaben ihre eigene Kompetenz und bestehende Verständnislücken blieben unentdeckt. Studierende, die keine Probleme bei der Bearbeitung hatte, konnten mit der KI schneller programmieren und erstellten damit Code, den sie ohnehin händisch ersetellen wollten. Außerdem konnten sie zwischen hilfreichen und unkorrekten Vorschlägen der KI unterscheiden. \textcite[S.~482-484]{prather_widening_2024} führen die Unterschiede darauf zurück, dass die Studierenden, die keine Schwierigkeiten hatten, KI-Vorschläge anhand eigener Lösungspläne kritisch hinterfragen konnten, während Lernende mit Schwierigkeiten sich vermehrt von den Vorschlägen der KI leiten ließen.

Auch experimentelle Untersuchungen außerhalb der Programmierausbildung zeigen, dass Aufgabenleistung und Kompetenzerwerb getrennt betrachtet werden müssen. \textcite[S.~1-3]{fan_beware_2025} untersuchten in einer randomisierten Studie die Bearbeitung einer Schreibaufgabe. Die Lernenden wurden in unterschiedliche Gruppen aufgeteilt, die unterschiedliche Hilfestellungen zur Erledigung der Aufgabe verwenden konnte. Die Gruppe, die ChatGPT als Unterstützung verwenden durfte, erzielte größere Verbesserungen in der Bewerteung des Texts, während sich beim Wissenszuwachs und Wissenstransfer keine statistisch signifikanten Unterschiede zu den anderen Gruppen zeigte. Die Studie von \textcite[S.~1]{bastani_generative_2025} liefert weitgehende Hinweise auf die Auswirkungen der Verwendung von generativer KI auf den Lernprozess.

Im randomisierten Feldexperiment wurden knapp 1000 High School Schüler*innen in zwei Gruppen aufgeteilt: \textit{GPT Base} und \textit{GPT Tutor}. Erstere erhielt ein klassisches KI-Interface, wie man es von ChatGPT kennt. Zweitere erhielt ein Interface mit Schutzmaßnahmen, die den Lernprozess unterstützen. Die GPT Base Gruppe erzielte zunächst bessere Übungsergebnisse, schnitten in einer anschließenden Prüfung ohne die Verwenung der Hilfsmittel 17\% schlechter ab als die Kontrollgruppe. Bei der didaktisch aufbereiteten Tutor-Variante wurde dieser negative Effekt weitestgehend verhindert. Die Studienautor*innen appellieren hierbei an die Entscheidungsträger, dass diese beim Einsatz von KI-Werkzeugen den langfristigen Wissenserwerb der Lernenden unbedingt beachten sollen \parencite[S.~1]{bastani_generative_2025}. Diese Befunde fanden zwar nicht im Feld der Informatik und im engeren Sinne der Programmierdidaktik statt, jedoch weisen beide Untersuchungen \textcite{fan_beware_2025} und \textcite{bastani_generative_2025} aber auch \textcite{prather_widening_2024} in eine Richtung und zwar, dass die Verwendung von generativer KI in Lehrbetrieb mit Bedacht eingesetzt werden soll und Lernende damit bessere Ergebnisse erzielen, als es ihnen ihr eigentlicher Kenntnisstand ermöglicht.

Einen wichtigen Beitrag für die Praxis liefern zudem spezialisierte Systeme wie CodeHelp von \textcite[S.~1]{CodeHelp} oder CodeAid von \textcite[S.~3]{KazemitabaarCodeAid}, die LLMs mit didaktischen Leitplanken kombinieren, so wie es auch \textcite[S.~1]{bastani_generative_2025} vorschlägt. Diese Tools verhindern, dass der User direkt eine Code-Lösung erhält. Sie bieten stattdessen erklärende Unterstützung, wodurch eigenständiges Denken angeregt werden soll. Bei \textcite[S.~9]{KazemitabaarCodeAid} beispielsweise werden die Leitplanken als Markierungen des fehlerhaften Code-Bereichs umgesetzt mit ergänzenden Hinweisen für die Studierenden. Erste Studien zu CodeAid und ähnlichen Classroom-Assistenten belegen, dass LLMs in realen Kursen bereits als standardisierte Unterstützungswerkzeuge eingesetzt werden können - allerdings nur, wenn klare Regeln zur Nutzung definiert werden. So wurde CodeHelp von \textcite{CodeHelp} in einer Folgeanalyse von \textcite[S.~1]{sheese_patterns_2024} an 52 Studierenden in einem data science Kurs angewandt. In den insgesamt 2500 untersuchten Anfragen an CodeHelp fanden die Autoren, dass die Studierenden, die das Tool am häufigsten verwendeten, auch besser im Kurs bewertet wurden. Zudem ergab die Untersuchung der Anfragen, dass am häufigsten Fragen zur dirketen Implementierung der Aufgabe oder dem Beheben eines Problems gestellt wurden. Fragen zu den Programmier-Konzepten der Aufgaben, die auf ein allgemeines Verstehen der Materie abzielen, wurden weniger gestellt \parencite[S.~54]{sheese_patterns_2024}.

Mit einem ähnlichen KI-Tutor namens \textit{Iris}, der ebenfalls didaktische Leitplanken aufweist und keine komplette Ergebnisse als Antwort gibt, untersuchen \textcite[S. ~1]{bassner_less_2026} die Bearbeitung einer Programmieraufgabe und den Wissenszuwachs danach. In der randomisierten Studie wurden die 275 Teilnehmer*innen in drei Gruppen aufgeteilt: (1) Gruppe \textit{Iris}, (2) ChatGPT und (3) herkömmliche Online-Ressourcen ohne KI-Unterstützung. In der Auswertung erreichte die Iris-Gruppe bei der Bewertung der Programmieraufgabe im Mittel 57,5 \% der Punkte und damit signifikant mehr als die Gruppe ohne KI-Unterstützung, die 29,8 \% der Punkte im Mittel erreichte. Die ChatGPT-Gruppe erreichte hierbei 71,8 \% also am meisten der möglichen Punkte aller Gruppen. Im Bereich des Wissenszuwachs zwischen Vor- und Nachtest sowie beim Verständnis des Codes zeigten sich keine statistisch signifikanten Unterschiede zwischen den drei Gruppen. Bei der Befragung gaben die Studierenden beider KI-Gruppen niedriegere Frustration und niedriegere kognitive Auslastung als die Studierenden der Gruppe ohne KI-Unterstützung an. Eine erhöhte intrinsische Motiviation konnte einzig bei der Iris-Gruppe festgestellt werden. Abschließend zogen \textcite[S.~15]{bassner_less_2026} den Schluss, dass KI in der Untersuchung ein \textit{Performance-Enhancer} ist, kein \textit{Learning-Enhancer}. Sie sprechen sich aufgrund der Vorteile bei der Motivation für die Verwendung für KI-Tools wie \textit{Iris} aus, da diese im Vergleich zu ChatGPT keine vollständigen Lösungen ausgeben und Studierende auf diesem Wege weniger eine sogennante \textit{comfort-trap} fallen würden. Denn die höhere Punkteanzahl bei der Aufgabe resultierte nicht in einem besseren Abschneiden beim Wissenstest danach.

Wie Studierende mit didaktischen Leitplanken umgehen, wenn sie diese eigenständig aus- oder einschalten können, untersuchten \textcite[S.~1]{kapoor_exploring_2026} mit dem Programmierassistenten der Lernumgebung Edugator. Die Funktion \textit{See Solution} erlaubt den teilnehmenden Studierenden bei der Erledigung der drei Debugging- und Programmieraufgaben die Leitplanken auszuschalten. Wenn ausgewählt, zeigt Edugator eine unregulierte Lösung der jeweiligen Aufgabe im ursprünglichen Wortlaut des Sprachmodells an. Die Autor*innen berichten von 50\% von den teilnehmenden 885 Studierenden, die die Option "See Solution" für mindestens eine Aufgabe verwendeten. 14 \% verwendeten es für alle Aufgaben. Vermehrt verwendet wurde diese Funktion von Studierenden, die schlecht bewertet wurden. Als Gründe gaben sie Hilfe beim Lösen des Problems, Zeitdruck und, am seltensten, fehlende Selbstregulation an \parencite[S.~ 29]{kapoor_exploring_2026}. 


\section{Generative KI in Prüfungen und Bewertungsszenarien}
Der Einsatz von LLMs stellt traditionelle Prüfungsformate grundlegend infrage. Studien wie die von \textcite[S.~6]{PintoLLMsForEducation} zeigen auf Lehrendenseite, dass LLMs offene Fragen in vielen Fällen ähnlich wie menschliche Bewertende beurteilen können - allerdings mit typischen Abweichungen: Manche Modelle sind strenger, andere übersehen subtile Fehler. Damit wird die Frage nach Fairness, Reliabilität und Transparenz solcher Bewertungsmethoden zentral.

Parallel dazu diskutieren \textcite[S.~1]{DennyPromptProblems2024}, \textcite[S.~1]{MahonGuidelines} und \textcite[S.~108]{PratherRobotsAreHere2023}, welche Aufgabenformate noch geeignet sind, individuelle Kompetenzen im LLM-Zeitalter zu prüfen. Aufgaben, die primär auf reiner Codeproduktion beruhen, verlieren an Aussagekraft, da LLMs solche Lösungen oft vollständig generieren. Stärker gefragt sind daher Problemanalyse, Codebewertung, Fehlersuche, Designentscheidungen und das Verfassen erklärender Texte. Diese Komponenten gelten als „LLM-resistent“, weil sie Verständnis und Reflexion voraussetzen und nicht vollständig automatisierbar sind.


\section{LLM-unterstütztes Programmierenlernen}
Mehrere kontrollierte Studien untersuchen, wie sich der Einsatz großer Sprachmodelle auf das Programmierlernen von Anfänger:innen auswirkt. Eine häufig zitierte Arbeit stammt von \textcite[S.~11]{KazemitabaarSupportLearners2023}: In einem randomisierten Experiment mit 69 Beginner:innen zeigte sich, dass Lernende mit Zugriff auf einen Codex-basierten Code-Generator signifikant häufiger funktionsfähige Programme erzeugten und ihre Aufgaben schneller lösten als die Kontrollgruppe. Gleichzeitig blieben ihre Leistungen bei nachfolgenden manuellen Code-Modifikationen stabil, was ein Hinweis darauf ist, dass der Einsatz von LLMs nicht automatisch zu Abhängigkeit führt.

Eine weitere Analyse von \textcite[S.~1]{kazemitabaar_how_2023} liefert Aufschluss über die Interaktionsmuster von Programmier-Einsteiger*innen bei der Verwendung von OpenAIs Codex. Anhand von Log-Daten identifizierten die Autoren vier distinkte Nutzungsstrategien bei Novizen: „AI Single Prompt“ (Generierung der gesamten Lösung), „AI Step-by-Step“, „Hybrid“ und „Manual“. Die Ergebnisse deuten darauf hin, dass die Art der Interaktion entscheidend für den nachhaltigen Lernerfolg ist: Der „Hybrid“-Ansatz, der eine Mischung aus eigenständigem Schreiben und gezielter KI-Nutzung darstellt, zeigte positive Trends hinsichtlich der Ergebnisse in späteren Retention-Tests. Wohingegen der „AI Single Prompt“-Ansatz eher mit negativen Resultaten korrelierte. Dies unterstreicht, dass der didaktische Mehrwert von LLMs maßgeblich vom Grad der Selbstregulation (Self-Regulation) und der eigenen kognitiven Beteiligung abhängt \parencite[S.~12]{kazemitabaar_how_2023}.

\textcite[S.~3-6]{XueIntro2024} untersuchten den Einsatz von ChatGPT 3.5 in einem CS1-Java-Kurs. Obwohl ChatGPT syntaktische Fehler reduzieren konnte, zeigten sich keine signifikanten Vorteile in tiefergehenden konzeptuellen Aufgaben. Die Autor:innen betonen, dass LLMs ohne didaktische Einbettung nur begrenzt Lernzuwächse erzeugen.
Ein weiteres Experiment von \textcite[S.~6-8]{TobaTests2024} demonstriert vor allem Effizienzgewinne: Mit ChatGPT lösten Studierende Programmieraufgaben etwa doppelt so schnell wie ohne. Die Codequalität blieb jedoch vergleichbar. Gleichzeitig weist die Studie auf Risiken wie missbräuchliche Nutzung in Prüfungssituationen hin.

Eine systematische Meta-Analyse von \textcite[S.~12-15]{Alanazi2025} bestätigt diese Befunde: KI-gestützte Lernsettings verbessern im Durchschnitt Performance und Bearbeitungszeit deutlich, zeigen aber keine konsistenten Effekte auf tiefes Konzeptverständnis.



\section{User-Experience-Messungen von LLM-Chatbots mittels User Experience Questionnaires (UEQ)}
Die Evaluierung von LLM-basierten Chatbots erfordert Messinstrumente, die über reine Leistungsmetriken hinausgehen und die subjektive Wahrnehmung der Interaktion erfassen. Als Standardinstrument hat sich hierbei der User Experience Questionnaire (UEQ) etabliert, da er sowohl pragmatische (zielgerichtete) als auch hedonische (nicht-zielgerichtete) Qualitätsaspekte abdeckt \parencite[S.~2-3]{Schrepp2023UEQHandbook}.

Aktuelle Studien bestätigen die Eignung dieses Messinstruments für Chatbots: So konnten \textcite[S.~17]{Mendoza2022} in einer Untersuchung im schulischen Umfeld nachweisen, dass der UEQ verlässliche Aussagen über die Akzeptanz von Lernassistenten liefert, wobei insbesondere die Dimensionen der Effizienz und der Verlässlichkeit (Efficieny und Dependability) als kritische Erfolgsfaktoren identifiziert wurden. \textcite[]{Ambacher2024} hebt in seiner Analyse von LLM-User-Interfaces hervor, dass eine mangelnde Transparenz über die Funktionsweise des Modells oder unerwartetes Systemverhalten direkt die Dimenson Verlässlichkeit negativ beeinflusst. Wird das System als unberechenbar wahrgenommen, sinkt die Akzeptanz, selbst wenn die generierten Antworten technisch hilfreich sein könnten. Demgegenüber verweist \textcite[S.~6]{Shahu2025} auf das Potential von LLMs, die pragmatische Qualität (Effizienz, Verständlichkeit und Verlässlichkeit) einer Aufgabe oder eines Produkts stark zu steigern. Ihre Arbeit handelt von einem GIS-Assistenten namens "MapGPT", der durch Implementierung natürlichsprachlicher Eingaben die Barrieren zur Nutzung dieser komplexer Software senken konnte. Dieser Sachverhalt schlug sich in positiven Werten auf den UEQ-Skalen für Efficiency (Effizienz) und Perspicuity (Verständlichkeit) nieder. \textcite[S.~5-6]{Lee2024} verwendeten auch den UEQ als Messinstrument für die Evaluierung ihres Chatbots "ChatFive" und konnten in der Dimension Attraktivität sowie dem Bereich Stimulation signifikant höhere Werte erzielen.

Folglich erweist sich der UEQ als das methodisch adäquate Instrument, um nicht nur die funktionale Gebrauchstauglichkeit (Pragmatik), sondern auch die durch generative KI neu entstehenden emotionalen Interaktionsqualitäten (Hedonik) in einem erprobten Messmodell vergleichbar zu erfassen.


\section{Zusammenfassung}
TODO
